\documentclass[11pt]{letter}
% \usepackage[utf8]{inputenc} % Required for inputting international characters
% \usepackage[T1]{fontenc} % Output font encoding for international characters

\usepackage{fouriernc} % Use the URW Schoolbook L font, similar to New Century Schoolbook
\usepackage{geometry} % Required for adjusting page dimensions and margins
\usepackage{hyperref}
\geometry{
paper=a4paper, % Paper size, change to letterpaper for US letter size
top=2cm, % Top margin
bottom=2.5cm, % Bottom margin
left=3.25cm, % Left margin
right=3.25cm, % Right margin
%showframe, % Uncomment to show how the type block is set on the page
}

\begin{document}
% \date{\raggedright\today} % Left-aligned date of the letter, use \today for the current date

\begin{letter}{
Dear Editors,
}

\begin{center}
\large\bfseries
Cover Letter for Submission -

\emph{Generalizable Dynamic Neural Ant Colony Search \\ for Large-Scale Routing Problems}
\end{center}

\opening{}

Please find enclosed our manuscript entitled ``\emph{Generalizable Dynamic Neural Ant Colony Search for Large-Scale Routing Problems}'' for consideration in \emph{IEEE Transactions on Artificial Intelligence}.
\vspace{6pt}

A preliminary version of this work appeared in the \emph{32nd ACM SIGKDD Conference on Knowledge Discovery and Data Mining (KDD-26)} under the title ``\emph{Beyond Static Priors: Dynamic Neural Guidance for Large-Scale Ant Colony Optimization}''. The conference paper introduced \textbf{Dynamic Neural Ant Colony Optimization (DYNACO)}, where the policy periodically observes the evolving pheromone distribution and incumbent solution, then updates edge-level guidance throughout the search trajectory. It also introduced the perturbation-based ACO backend and scope-restricted refinement procedure that make dynamic neural guidance practical for large-scale routing problems.

This journal submission substantially extends the conference version with \textbf{Dynamic Neural Ant Colony Search (DYNACS)}, which makes the dynamic neural colony population-aware. The conference version made ACO guidance state-aware; the journal version lets different ant groups express complementary behavior-conditioned trajectory distributions while sharing one incumbent solution, pheromone field, and refinement process. The main journal extensions are summarized as follows.
\vspace{-5pt}

\begin{enumerate}
\item We introduce \textbf{DYNACS}, which extends DYNACO from one state-aware transition bias to \textbf{trajectory-diverse neural ACO}:
\begin{enumerate}
\item dynamic guidance is parameterized as a \textbf{mixture of behavior-conditioned trajectory distributions}, allowing different ant groups to express complementary search behaviors within one shared ACO process;
\item multiple behavior-conditioned guidance fields are generated from the same evolving ACO macro-state, supporting heterogeneous exploration and exploitation patterns across the ant population;
\item all behaviors share the same incumbent solution, pheromone field, and refinement process, so successful perturbations from one behavior immediately shape the common search memory used by the others.
\end{enumerate}

\item We add \textbf{learned behavior allocation} for DYNACS:
\begin{enumerate}
    \item a state-conditioned allocator assigns ant mass among learned behaviors from the encoded ACO macro-state;
    \item the allocation policy works inside the shared colony state, preserving a single pheromone process and a single local-refinement backend;
    \item the resulting colony adapts the exploration--exploitation balance while retaining the efficiency and reliability of the original DYNACO search backbone.
\end{enumerate}

\item We substantially expand the experimental evaluation under a stricter \textbf{zero-shot benchmark generalization} protocol:
\begin{enumerate}
    \item a single model trained only on uniformly distributed 1K-node synthetic instances is directly tested on diverse benchmark TSP and CVRP instances with different scales, structures, and distributions, without retraining, fine-tuning, or instance-specific tuning;
    \item DYNACS achieves an average gap of \textbf{1.22\%} on TSP and \textbf{2.80\%} on CVRP across benchmark datasets, outperforming state-of-the-art neural baselines under matched evaluation;
    \item compared with DYNACO, DYNACS improves the total average gap from \textbf{1.48\%} to \textbf{1.22\%} on TSP and from \textbf{3.10\%} to \textbf{2.80\%} on CVRP, connecting the journal extension to population-aware guidance under the same dynamic ACO search process.
\end{enumerate}

\item We provide new analyses and ablation studies for the journal extension:
\begin{enumerate}
    \item the ablations compare the unguided ACO backend, the state-aware DYNACO model, and the population-aware DYNACS model under a matched search budget;
    \item the studies connect the improvement to behavior-conditioned trajectory guidance, shared search memory, and learned state-conditioned allocation;
    \item guidance analyses show how the learned policy adapts to the evolving pheromone landscape, maintains complementary search behaviors, and strengthens zero-shot benchmark generalization.
\end{enumerate}


\end{enumerate}
\vspace{6pt}

We believe that this work makes a significant contribution to learning-guided optimization, machine learning for combinatorial optimization, and neural-guided metaheuristics. The journal extension shows that dynamic neural guidance can improve iterative solvers most effectively when it both adapts to the solver state and preserves useful behavioral diversity throughout the search trajectory. Our source code, datasets, and pretrained model checkpoints are publicly available at \href{https://doi.org/10.5281/zenodo.20502184}{https://doi.org/10.5281/zenodo.20502184}, and the implementation repository is available at \href{https://github.com/shoraaa/DyNACO}{https://github.com/shoraaa/DyNACO}. Thank you for considering our submission.
\vspace{6pt}

\closing{Sincerely, \\ The authors}

\end{letter}

\end{document}
